\documentclass[11pt,a4paper]{altacv}

\geometry{left=1.5cm,right=1.5cm,top=1.cm,bottom=1.2cm,columnsep=1.5cm}

\usepackage[utf8]{inputenc}
\usepackage[T1]{fontenc}
\usepackage[french]{babel}
\usepackage{paracol}
\usepackage{xcolor}
\usepackage{hyperref}
\usepackage{fontawesome5}
\usepackage{setspace}

\definecolor{SlateGrey}{HTML}{2E2E2E}
\definecolor{LightGrey}{HTML}{666666}
\definecolor{DarkPastelRed}{HTML}{8AC0D9}
\definecolor{PastelRed}{HTML}{00B0DA}
\definecolor{GoldenEarth}{HTML}{B4AD0C}
\colorlet{name}{black}
\colorlet{tagline}{PastelRed}
\colorlet{heading}{DarkPastelRed}
\colorlet{headingrule}{GoldenEarth}
\colorlet{subheading}{PastelRed}
\colorlet{accent}{PastelRed}
\colorlet{emphasis}{SlateGrey}
\colorlet{body}{LightGrey}

\renewcommand{\familydefault}{\sfdefault}

\name{\Large Guillaume BOILEAU}
\tagline{Chef de projet technique SI | SPOC, besoins techniques, build/run, risques, cycle en V \& coordination DSI}
\personalinfo{
  \email{guillaume.boileaupro@gmail.com}
  \phone{+33 6 59 94 85 84}
  \location{Alpes-Maritimes, France}
  \linkedin{guillaume-boileau-phd-891b81265}
}

\setlength{\parskip}{0.72\baselineskip}

\begin{document}
\makecvheader
\vspace{-0.65cm}

\cvsection{Lettre de motivation}
\vspace{-0.45cm}

{\color{accent}\textbf{Objet :}} \textit{Candidature -- Chef de projet stratégique informatique -- Valbonne}

{\color{body}

Madame, Monsieur,

Je vous adresse ma candidature pour le poste de \textbf{Chef de projet stratégique informatique} au sein de la Direction Technique de la CNAF. Votre offre m'intéresse particulièrement car elle décrit un rôle de coordination technique transverse au cœur d'une DSI : être le point d'entrée des demandes techniques, qualifier les besoins, consolider les contraintes, représenter les impacts build/run, fédérer les parties prenantes, identifier les risques, préparer les arbitrages et assurer un reporting clair auprès des différents acteurs.

Mon parcours s'est construit dans des environnements techniques exigeants, à l'interface entre \textbf{développement logiciel, systèmes complexes, validation, support opérationnel, documentation et coordination}. Je ne viens pas d'un parcours purement administratif de chefferie de projet : mon expérience s'est construite au contact des systèmes, des données, des contraintes techniques, des anomalies et des livrables. C'est précisément ce qui me permet de comprendre un besoin technique, de poser les bonnes questions, de repérer les dépendances, de rendre visibles les risques et de transformer un sujet complexe en plan d'action lisible.

Au \textbf{CNES}, j'ai travaillé sur le segment sol de la mission spatiale \textbf{LISA}, dans un environnement CNES/ESA associant chaînes de données mission, niveaux L0--L3, cycle en V, intégration de modèles, validation technique et revues collaboratives. J'y ai développé \textbf{CATpyLISA}, une plateforme Python destinée à structurer des workflows, intégrer des contributions hétérogènes, comparer des résultats, qualifier des écarts et documenter les décisions techniques. Cette expérience m'a donné une pratique concrète du cadrage de besoins, de la traçabilité, de la validation, de l'analyse d'écarts, de la préparation de revues et du transfert d'information vers plusieurs contributeurs.

Dans ce cadre, il fallait gérer des sujets très proches de ceux décrits dans votre offre : clarifier les besoins techniques, identifier les hypothèses, suivre les dépendances, documenter les contraintes, produire des synthèses, faire remonter les points bloquants et maintenir un niveau d'information partagé entre acteurs. La logique du cycle en V, la structuration des livrables, l'analyse des risques techniques et la validation des résultats faisaient partie intégrante de mon travail. Cette expérience me semble directement transposable au rôle attendu pour la Direction Technique : qualifier les entrants, préparer les éléments nécessaires au lancement des projets techniques et assurer une continuité claire entre demande, instruction, décision et transfert vers le chef de projet désigné.

Plus récemment, chez \textbf{VEV Platform Services France}, j'ai travaillé dans un environnement logiciel industriel connecté à des infrastructures terrain. J'y ai contribué à l'intégration de services applicatifs, à l'analyse de flux de données, au support opérationnel, à l'investigation d'incidents et à la documentation technique. Cette expérience m'a permis de renforcer ma compréhension des enjeux \textbf{build/run} : continuité de service, analyse de comportements applicatifs, flux opérationnels, diagnostic, correction, suivi des anomalies et qualité de service. Elle m'a également appris à formaliser clairement des constats techniques pour faciliter la résolution par les équipes concernées.

Le poste proposé par la CNAF demande également de travailler avec plusieurs directions et de fédérer des parties prenantes variées : Direction Technique, Direction des Opérations, chefs de projets d'autres directions, architectes techniques, équipes de production, support et performance du run. C'est un mode de fonctionnement dans lequel je peux apporter de la valeur. J'ai l'habitude de dialoguer avec des profils différents, de traduire des contraintes techniques en éléments compréhensibles, de structurer l'information et de garder une communication constructive même lorsque les sujets sont complexes ou bloquants.

Je peux notamment contribuer sur les missions suivantes : récupérer et qualifier les besoins techniques, identifier les contraintes de planning et les dépendances, consolider les impacts techniques, suivre les actions ouvertes, préparer des synthèses pour les comités, documenter les risques, proposer des plans d'action et assurer un passage de relais propre vers les équipes chargées du pilotage opérationnel. Mon expérience ne se limite pas à produire de l'analyse : elle m'a habitué à rendre les sujets exploitables, suivables et décidables.

Je dispose aussi d'un socle technique utile pour dialoguer efficacement avec les équipes de la DSI : Python, TypeScript, Angular, Node.js, Linux, Git, Docker, CI/CD, SQL, monitoring/logs, OpenSearch, Sentry, Confluence, Jira et Zenhub. Je ne me présente pas comme expert infrastructure DSI pur, mais comme un profil capable de comprendre les enjeux techniques, de les qualifier correctement, de faciliter la coordination entre spécialistes et de maintenir une vision claire des impacts sur le build, le run, la validation et la qualité de service.

La mission de la CNAF donne enfin un sens particulier à ma candidature. Contribuer à une plateforme technique au service des équipes de développement, dans une DSI portant des services publics essentiels, correspond à ce que je recherche aujourd'hui : un poste transverse, exigeant, utile, où la rigueur technique, la qualité de coordination et la clarté de l'information ont un impact concret sur le fonctionnement collectif.

Je serais heureux de pouvoir échanger avec vous afin de vous présenter plus en détail mon parcours et la manière dont mon expérience en coordination technique, cycle en V, validation, support run, documentation et analyse de risques peut contribuer aux missions de la Direction Technique de la CNAF.

Je vous prie d'agréer, Madame, Monsieur, l'expression de mes salutations distinguées.

\textbf{Guillaume Boileau}

}

\end{document}
